\section{Additionner des nombres relatifs}

\begin{propbox}
\bb{Règle 1 -- Même signe}

Pour additionner deux nombres relatifs de même signe :

\begin{itemize}
    \item on additionne les distances à zéro ;
    \item on conserve le signe commun.
\end{itemize}
\end{propbox}

\begin{exemplebox}
\[
(+5)+(+7)=+12
\]

\[
(-5)+(-10)=-15
\]
\end{exemplebox}

\begin{propbox}
\bb{Règle 2 -- Signes contraires}

Pour additionner deux nombres relatifs de signes contraires :

\begin{itemize}
    \item on soustrait les distances à zéro ;
    \item on garde le signe du nombre qui a la plus grande distance à zéro.
\end{itemize}
\end{propbox}

\begin{exemplebox}
\[
(+5)+(-7)=-2
\]

\[
(-5)+(+10)=+5
\]
\end{exemplebox}

\begin{methodebox}
Pour calculer une somme :

\begin{enumerate}
    \item Je regarde les signes.
    \item Je choisis la bonne règle.
    \item Je calcule les distances à zéro.
    \item Je donne le bon signe.
\end{enumerate}
\end{methodebox}
